% page 223 du fac-simile (image p-222 du scan)
% bloc 188.0 x 250.6 mm ; marge gauche 14.0 mm
\pgc{-7.620mm}{0.000mm}{
\l{}
\l{}
\l{\cel{63}215}
\l{}
\l{\cel{6}\sou{\textquotedbl{}hidalgo\textquotedbl{}}. Nobelo Hispana qua asertas decendar de familio Krist-}
\l{\cel{9}ana anciena. - DEFIRS.}
\l{}
\l{\cel{6}hidatodo. (\sou{biol.}) Nomo kreita da Haberlandt por signifikar organo}
\l{\cel{9}mikroskopala qua ekpulsas la aquo quan kontenas la planto.}
\l{}
\l{\cel{6}hidro. I. (\sou{mitol}.)Serpento sep-kapa, qua produktis sep altra}
\l{\cel{9}kapi kande onu tranchis un ek la sep. - II. (\sou{zool.)}\cel{2}Aquo-ser-}
\l{\cel{9}pento reptera ; anke : mikra polipo qua habitas aquo sen-sala e}
\l{\cel{9}su nutras ek mikra krustacei e dafni, quin lu kaptas per sua}
\l{\cel{9}palpili. - L.\cel{1}\sou{hydra}. - DEFIRS.}
\l{}
\l{\cel{6}\sou{hidramnio}to. (\sou{anat.}) Hidropso di la amniono ; hiper-abundo di la}
\l{\cel{9}liquido amnionala.}
\l{}
\l{\cel{6}\sou{hidranto}. Buxo ek gisfero, qua kontenas junto-tubo e robineto}
\l{\cel{9}por aquo-cherpo, instalita sive borde di la choseo, sive an}
\l{\cel{9}muro, ed uzata da la extingisti di incendii. - DE.\cel{2}s\cel{2}q c i}
\l{}
\l{\cel{6}\sou{hidrato}. (\sou{kemio}). Kompozajo derivita ek aquo, per substituco, o}
\l{\cel{9}formacita per la kombino nemediata di substanco kun aquo.- DEFIS.}
\l{}
\l{\cel{6}\sou{hidrauliko}. La parto di fiziko qua traktas la +leyi (legi) pri}
\l{\cel{9}la movo da la liquidi.\cel{2}- DEFIRS.}
\l{}
\l{\cel{6}\sou{hidrocefala}. (\sou{anat.}) Hidropso di la kapo. - DEF.}
\l{}
\l{\cel{6}\sou{hidrodinamiko}. La parto di hidrauliko qua traktas la movi e la}
\l{\cel{9}presi da la fluidi. - DEFIRS.}
\l{}
\l{\cel{6}\sou{hidrofila}. I.Qua esas avida ye aquo, qua absorbas aquo. II. (\sou{Mate}-}
\l{\cel{9}\sou{rio koloida}) qua inflesas en aquo.}
\l{}
\l{\cel{6}\sou{hidrofobio}. (\sou{patol}.)Hororo di la liquidi, qua karakterizas ula}
\l{\cel{9}morbi (kom ex.: rabio). - DEFIRS.}
\l{}
\l{\cel{6}\sou{hidrogeno}. (\sou{kemio)}. Sinonimo di \textquotedbl{}hido\textquotedbl{}. - DEFIRS.}
\l{}
\l{\cel{6}\sou{hidrografio}. Cienco qua havas kom temo la studio e la deskripto}
\l{\cel{9}di la riveri, fluvii e mari. - DEFIS.}
\l{}
\l{\cel{6}\sou{hidrografo}. La persono qua esas kompetenta pri hidrografio.}
\l{}
\l{\cel{6}\sou{hidroleucito.}\cel{2}(\sou{biol}.) Vakuolo kontenata\cel{2}en la celuli di la pro-}
\l{\cel{9}plasmo di la vejetanti, provizita ye membrano propra (+tono-}
\l{\cel{9}plasto) e qua su reproduktas per dupartigo.}
\l{}
\l{\cel{6}hidrolizo. Reakto di duimesko di molekulo per la efiko da aquo. -}
\l{\cel{9}DEFIS.}
\l{}
\l{\cel{6}hidrologio. Parto di la natur-historio qua traktas aquo, e, preci-}
\l{\cel{9}pue, la aqui mineraloza. - DEFIS.}
\l{}
\l{\cel{6}hidromancio.\cel{2}Divinado per aquo. - DEFIS}
}
